%=================================%
\section{Extended Parameter Space}
\label{app:parameter_space}
%=================================%

In Section \ref{sec:eccentric_torques} we focused our discussion on the  $h_\textrm{p} = 0.06$ discs, while varying other disc parameters. This is the lowest value of $h_\textrm{p}$ in our main parameter sweep and therefore allows us to probe the largest values of the ratio $\tilde{e} = e/h_\textrm{p}$. We also computed the torque and migration timescale results for higher values of $h_\textrm{p} = 0.08$ and $0.1$, allowing us to explore the variation of torque with aspect ratio. In Figs.~\ref{fig:torque_eccentricity_h_0.08} and \ref{fig:torque_eccentricity_h_0.1} we plot the different torques as a function of eccentricity across $(q,p)$ space, akin to Fig.~\ref{fig:torque_eccentricity_h_0.06}. Then in Figs.~\ref{fig:eccentric_timescales_h_0.08} and \ref{fig:eccentric_timescales_h_0.1} we do the same for the migration rates, as per the fiducial Fig.~\ref{fig:eccentric_timescales_h_0.06} in the main text. 

Furthermore, we make our full data cube available to the community as supplementary material to the journal and also as a public repository online (\href{https://github.com/callumf8/Eccentric-planet-disc-interactions.git}{https://github.com/callumf8/Eccentric-planet-disc-interactions.git}). Accessing and manipulating this data cube is described more in the accompanying \texttt{README.md} file. Furthermore, we have provided an exemplar jupyter notebook script \texttt{read\_data.ipynb}, which performs the reading in of the $(q,p,h_\mathrm{p})$ data cube, interpolation onto intermediate regions of parameter space and demonstrates how to calculate migration tracks as shown in Fig.~\ref{fig:1d_migration}. Whilst some extrapolation outwith the bounds of our study is perhaps justified, one should be careful not to extend it too far as our results demonstrate the complicated and sensitive dependence on the underlying disc parameters and planetary eccentricities.
%
\begin{figure*}
    \centering
    \includegraphics[width=0.9\textwidth]{Figures/torque_eccentricity_h_0.08.pdf}
    \caption{As per Fig.~\ref{fig:torque_eccentricity_h_0.06} but for $h_\textrm{p}=0.08$.}
    \label{fig:torque_eccentricity_h_0.08}
\end{figure*}
%
\begin{figure*}
    \centering
    \includegraphics[width=0.9\textwidth]{Figures/torque_eccentricity_h_0.1.pdf}
    \caption{As per Fig.~\ref{fig:torque_eccentricity_h_0.06} but for $h_\textrm{p}=0.10$.}
    \label{fig:torque_eccentricity_h_0.1}
\end{figure*}
%
\begin{figure*}
    \centering
    \includegraphics[width=0.9\textwidth]{Figures/eccentric_timescales_h_0.08.pdf}
    \caption{As per Fig.~\ref{fig:eccentric_timescales_h_0.06} but for $h_\textrm{p}=0.08$.}
    \label{fig:eccentric_timescales_h_0.08}
\end{figure*}
%
\begin{figure*}
    \centering
    \includegraphics[width=0.9\textwidth]{Figures/eccentric_timescales_h_0.1.pdf}
    \caption{As per Fig.~\ref{fig:eccentric_timescales_h_0.06} but for $h_\textrm{p}=0.10$}
    \label{fig:eccentric_timescales_h_0.1}
\end{figure*}

